\section*{Chapitre \ref{ch:g11:organic-skeletons} --- Les molécules organiques~: squelettes et noms}

\begin{solution}{exo:g11:organic-skeletons:1}
Les deux ont six \omterm{def:g7:atoms-and-molecules:atom}{atomes} de carbone (extrémités et sommets) et la formule
\ce{C6H14}~: l'hexane et le 2-méthylpentane sont \omterm{def:g11:organic-skeletons:isomer}{isomères}.
\end{solution}

\begin{solution}{exo:g11:organic-skeletons:2}
Propane~; hexane~; 2-méthylbutane.
\end{solution}

\begin{solution}{exo:g11:organic-skeletons:3}
\ce{CH3-CH2-CH3}; \chemfig{CH_3-CH(-[2]CH_3)-CH_3};
\chemfig{H-C(-[2]H)(-[6]H)-C(-[2]H)(-[6]H)-H}.
\end{solution}

\begin{solution}{exo:g11:organic-skeletons:4}
La première, le 3-méthylpentane~: 6 \omterm{def:g7:atoms-and-molecules:atom}{atomes} de carbone~; ses extrémités
\ce{CH3} (il y en a trois) portent 9 \omterm{def:g7:atoms-and-molecules:atom}{atomes} d'hydrogène, ses deux sommets
\ce{CH2} en portent 4 et son point de ramification \ce{CH} 1~:
\ce{C6H14}. Le cyclohexane~: 6 sommets, chacun \ce{CH2}~: \ce{C6H12}.
\end{solution}

\begin{solution}{exo:g11:organic-skeletons:5}
\ce{C8H18}. $2n + 2 = 20$ donne $n = 9$~: le nonane, \ce{C9H20}.
\end{solution}

\begin{solution}{exo:g11:organic-skeletons:6}
\setchemfig{atom sep=1.6em}
3-méthylhexane \chemfig{-[:30]-[:-30]-[:30](-[:90])-[:-30]-[:30]-[:-30]},
\ce{C7H16}~; 2,2-diméthylbutane \chemfig{-[:0](-[:90])(-[:-90])-[:0]-[:30]},
\ce{C6H14}~; 3-éthylpentane
\chemfig{-[:30]-[:-30](-[:-90]-[:-30])-[:30]-[:-30]}, \ce{C7H16}.
\end{solution}

\begin{solution}{exo:g11:organic-skeletons:7}
Le pentane, le 2-méthylbutane et le 2,2-diméthylpropane, dessinés dans la
figure du chapitre.
\end{solution}

\begin{solution}{exo:g11:organic-skeletons:8}
Dans le «~2-éthylbutane~», la plus longue chaîne compte cinq carbones en
passant par le groupe éthyle~: c'est le 3-méthylpentane. Le
«~4-méthylpentane~» est numéroté à partir de la mauvaise extrémité~:
c'est le 2-méthylpentane. Le «~1-méthylbutane~» est simplement une
chaîne de cinq carbones~: c'est le pentane.
\end{solution}

\begin{solution}{exo:g11:organic-skeletons:9}
Le 2,2-diméthylpropane (\qty{9.5}{\celsius}) < le 2-méthylbutane
(\qty{27.8}{\celsius}) < le pentane (\qty{36.1}{\celsius}). Les mêmes
\omterm{def:g7:atoms-and-molecules:atom}{atomes} sont disposés de façon de plus en plus compacte~: moins de contact
entre \omterm{def:g7:atoms-and-molecules:molecule}{molécules}, des \omterm{def:g11:polarity-and-cohesion:van-der-waals}{interactions de van der Waals} plus faibles, une
température d'ébullition plus basse.
\end{solution}

\begin{solution}{exo:g11:organic-skeletons:10}
Refermer la chaîne en cycle crée une liaison \ce{C-C} de plus, qui prend
la place de deux liaisons \ce{C-H} (une à chaque extrémité). Les
\omterm{def:g11:organic-skeletons:formulas}{formules brutes} sont différentes~: ce ne sont pas des \omterm{def:g11:organic-skeletons:isomer}{isomères}. Le
cyclohexane n'est pas un \omterm{def:g11:organic-skeletons:alkane}{alcane} (son squelette est un cycle~; sa formule
n'est pas \ce{C_{n}H_{2n+2}}).
\end{solution}

\begin{solution}{exo:g11:organic-skeletons:11}
Le butane et le 2-méthylpropane (\ce{C4H10}). Le propane a pour formule
\ce{C3H8} et le cyclobutane \ce{C4H8}.
\end{solution}

\begin{solution}{exo:g11:organic-skeletons:12}
\setchemfig{atom sep=1.5em}
\begin{center}
\small
\begin{tabular}{ccc}
\chemfig{-[:30]-[:-30]-[:30]-[:-30]-[:30]} &
\chemfig{-[:30](-[:90])-[:-30]-[:30]-[:-30]} &
\chemfig{-[:30]-[:-30](-[:-90])-[:30]-[:-30]} \\
hexane & 2-méthylpentane & 3-méthylpentane \\[6pt]
\chemfig{-[:0](-[:90])(-[:-90])-[:0]-[:30]} &
\chemfig{-[:30](-[:90])-[:-30](-[:-90])-[:30]} & \\
2,2-diméthylbutane & 2,3-diméthylbutane &
\end{tabular}
\end{center}
\end{solution}

\begin{solution}{exo:g11:organic-skeletons:13}
Deux longs zigzags peuvent s'aligner côte à côte sur toute leur longueur,
avec beaucoup d'\omterm{def:g7:atoms-and-molecules:atom}{atomes} en contact~; deux boules ne se touchent que sur une
petite surface. Les \omterm{def:g11:polarity-and-cohesion:van-der-waals}{interactions de van der Waals} agissent entre \omterm{def:g7:atoms-and-molecules:atom}{atomes}
en contact~: elles sont plus grandes entre \omterm{def:g7:atoms-and-molecules:molecule}{molécules} linéaires, qu'il
faut chauffer davantage pour les séparer.
\end{solution}

\begin{solution}{exo:g11:organic-skeletons:14}
La plus longue chaîne compte six carbones~; numérotés par la gauche, les
substituants sont en 2, 4, 4, par la droite en 3, 3, 5~; la première
différence (2 contre 3) décide~: 2,4,4-triméthylhexane.
\end{solution}

\begin{solution}{exo:g11:organic-skeletons:15}
\setchemfig{atom sep=1.6em}
\chemfig{-[:0](-[:90])(-[:-90])-[:-30]-[:30](-[:90])-[:-30]}~: une chaîne
de cinq carbones avec deux groupes méthyle sur le carbone 2 et un sur le
carbone 4, soit huit \omterm{def:g7:atoms-and-molecules:atom}{atomes} de carbone en tout~; $2 \times 8 + 2 = 18$
\omterm{def:g7:atoms-and-molecules:atom}{atomes} d'hydrogène, \ce{C8H18}, la formule de l'octane~: ce sont des
\omterm{def:g11:organic-skeletons:isomer}{isomères}. \omterm{def:g11:organic-skeletons:formulas}{Formule semi-développée}~:
\chemfig{CH_3-C(-[2]CH_3)(-[6]CH_3)-CH_2-CH(-[2]CH_3)-CH_3}.
\end{solution}

\begin{solution}{pb:g11:organic-skeletons:1}
\setchemfig{atom sep=1.5em}
\textbf{1.} \ce{C4H10}~: avec $n = 4$, $2n + 2 = 10$.

\textbf{2.} \chemfig{H-C(-[2]H)(-[6]H)-C(-[2]H)(-[6]H)-C(-[2]H)(-[6]H)-C(-[2]H)(-[6]H)-H}.

\textbf{3.} \ce{CH3-CH2-CH2-CH3}~; \chemfig{-[:30]-[:-30]-[:30]}.

\textbf{4.} \chemfig{H-C(-[2]H)(-[6]H)-C(-[6]H)(-[2,1.5]C(-[1]H)(-[3]H)(-[2]H))-C(-[2]H)(-[6]H)-H}~;
\chemfig{CH_3-CH(-[2]CH_3)-CH_3}~; \chemfig{-[:30](-[:90])-[:-30]}. Il
possède les mêmes \omterm{def:g7:atoms-and-molecules:atom}{atomes}, quatre carbones et dix hydrogènes, liés dans
un ordre différent.

\textbf{5.} Oui~: des \omterm{def:g11:organic-skeletons:isomer}{isomères de constitution} (squelettes différents).

\textbf{6.} Tous les trois~: leurs températures d'ébullition sont
inférieures à \qty{20}{\celsius}.

\textbf{7.} À \qty{-5}{\celsius}, au-dessous de sa température
d'ébullition, le butane reste liquide sous la pression normale et ne
dégage que peu de gaz.

\textbf{8.} Le propane (\qty{-42}{\celsius}) et le 2-méthylpropane
(\qty{-11.7}{\celsius}) bouillent au-dessous des températures
hivernales~: ils donnent encore beaucoup de gaz dans le froid.

\textbf{9.} Le 2-méthylpropane est ramifié, plus compact~: ses
\omterm{def:g11:polarity-and-cohesion:van-der-waals}{interactions de van der Waals} sont plus faibles.

\textbf{10.} Avec trois \omterm{def:g7:atoms-and-molecules:atom}{atomes} de carbone, le seul squelette possible est
la chaîne \ce{C-C-C}~: déplacer un carbone ailleurs redonne la même
chaîne.

\textbf{11.} Le 2-méthylbutane.

\textbf{12.} La chaîne doit être numérotée à partir de l'extrémité la
plus proche du substituant~: 2, et non 3.

\textbf{13.} 2, 3 et 5.

\textbf{14.} L'heptane, \chemfig{-[:30]-[:-30]-[:30]-[:-30]-[:30]-[:-30]}.

\textbf{15.} Le 2-méthylhexane \chemfig{-[:30](-[:90])-[:-30]-[:30]-[:-30]-[:30]}
et le 3-méthylhexane \chemfig{-[:30]-[:-30](-[:-90])-[:30]-[:-30]-[:30]}.
Le nom «~4-méthylhexane~» désigne en réalité le 3-méthylhexane,
numéroté à partir de la mauvaise extrémité.

\textbf{16.} Le 2,2-, le 2,3-, le 2,4- et le
3,3-diméthylpentane.

\textbf{17.} Oui, le 3-éthylpentane. Sur le carbone 2, le groupe éthyle
formerait une chaîne de six carbones passant par lui~: la \omterm{def:g7:atoms-and-molecules:molecule}{molécule} serait
le 3-méthylhexane.

\textbf{18.} Le 2,2,3-triméthylbutane,
\chemfig{-[:0](-[:90])(-[:-90])-[:0](-[:90])-[:0]}.

\textbf{19.} $1 + 2 + 4 + 1 + 1 = 9$.

\textbf{20.} L'heptane possède 9 \omterm{def:g11:organic-skeletons:isomer}{isomères de constitution}, lui-même
compris.
\end{solution}
